% =========================== Fiche de tache, resume et abreviations =========
\chapter*{Fiche de tâche et résumé}
\markboth{Fiche de tâche et résumé}{}
\addcontentsline{toc}{chapter}{Fiche de tâche et résumé}

\begin{table}[H]
\centering
\small
\begin{tabular}{@{}R{0.26}R{0.74}@{}}
\hautfilet
\textbf{Rubrique} & \textbf{Contenu}\\
\medfilet
Identifiant & T005\\
Intitulé & Mettre en place le DNS\\
Responsable & FOMETHE SOBMBANANG MAXIMILIEN\\
Pondération et priorité & 8 sur 10, priorité haute\\
Cellule & Réseaux et Sécurité, lot Services Network\\
Prérequis consommés & T001, plan d'adressage et modèle IPAM figés par le
  document 02 ; T003, réseau de management\\
Tâches coordonnées & T004 DHCP haute disponibilité, T006 synchronisation
  temporelle, T007 agent réseau et API de provisionnement\\
Livrable & Le présent dossier : conception, configurations, procédures et
  critères de recette du service DNS\\
Hors périmètre & T009, tests, validation, reprise et documentation, qui
  exploite les critères définis au chapitre~\ref{ch:recette} une fois les
  autres lots disponibles\\
\basfilet
\end{tabular}
\end{table}

\section*{Résumé}

Ce dossier définit et paramètre le service de résolution de noms de GANDAL. Il
sépare strictement le rôle autoritatif du rôle récursif, retient PowerDNS
conformément à l'arbitrage du document 03, et construit l'espace de nommage
interne à partir du plan d'adressage arrêté par le document 02. Il traite la
confidentialité des noms entre domaines par une politique de sélection des
zones indexée sur l'adresse source, doublée d'un filtrage réseau qui interdit
aux VNets tenant d'atteindre les serveurs autoritatifs.

Le cycle de vie des enregistrements est confié à un propriétaire d'écriture
unique, l'agent réseau, au moyen de l'API autoritative exposée derrière une
terminaison TLS. Les opérations sont rejouables sans effet de bord, et une
panne de publication laisse l'opération dans un état dégradé observable sans
libérer une adresse encore utilisée.

Les configurations proposées sont complètes et cohérentes avec les contraintes
matérielles du projet : une carte réseau par hôte, un commutateur unique, une
MTU utile de 1\,450 octets imposée par l'encapsulation VXLAN. Conformément
à la règle du corpus, aucun déploiement ni aucun essai n'est présenté comme
effectué : les contrôles de recette sont énoncés avec leur commande, leur
résultat attendu et le statut \statut{non exécuté}.

\begin{encadre}
\titreencadre{Mots-clés}
GANDAL ; DNS ; PowerDNS ; zone autoritative ; résolveur récursif ; espace de
nommage interne ; isolation multi-tenant ; TSIG ; DNSSEC ; publication
idempotente ; mode dégradé.
\end{encadre}

\section*{Abréviations propres au dossier}

Le glossaire général du corpus reste applicable. Les termes ci-dessous sont
ceux que ce dossier emploie de manière spécifique.

\begin{table}[H]
\centering
\small
\begin{tabular}{@{}R{0.17}R{0.83}@{}}
\hautfilet
\textbf{Terme} & \textbf{Définition retenue}\\
\medfilet
Autoritatif & Serveur qui détient les données d'une zone et répond avec
  l'indicateur \opt{AA}. Il ne résout pas les noms des autres zones.\\
Récursif & Résolveur qui interroge la chaîne des serveurs autoritatifs pour le
  compte d'un client et met les réponses en cache.\\
AXFR, IXFR & Transfert complet ou incrémental d'une zone entre serveurs.\\
TSIG & Signature symétrique des messages DNS, définie par le RFC 8945,
  employée ici pour les transferts de zone.\\
NTA & \textit{Negative Trust Anchor} : exception locale de validation DNSSEC
  sur une branche de l'arbre.\\
RRset & Ensemble des enregistrements partageant le même nom, la même classe et
  le même type ; unité de modification de l'API.\\
PTR & Enregistrement de résolution inverse, d'une adresse vers un nom.\\
Vue & Ensemble de zones qu'une catégorie de clients est autorisée à résoudre ;
  déterminée ici par l'adresse source.\\
dnstap & Format et transport de journalisation structurée des échanges DNS.\\
EDNS & Extension du protocole DNS permettant notamment d'annoncer une taille de
  réponse UDP supérieure à 512 octets.\\
\basfilet
\end{tabular}
\end{table}

\section*{Notations}

Le dossier emploie les notations suivantes, dont les définitions complètes
figurent aux sections~\ref{sec:nommage} et~\ref{sec:dimensionnement}.

\begin{table}[H]
\centering
\small
\begin{tabular}{@{}R{0.13}R{0.87}@{}}
\hautfilet
\textbf{Symbole} & \textbf{Signification}\\
\medfilet
$n$ & Numéro d'un tenant, $1\leqslant n\leqslant T$\\
$T$ & Nombre de tenants du périmètre minimal, $T=50$\\
$\mathcal{P}(n)$ & Préfixe IPv4 attribué au tenant $n$\\
$\mathcal{D}(n)$, $\mathcal{R}(n)$ & Zone directe et zone inverse du tenant $n$\\
$\tau(a)$ & Étiquette de vue associée à l'adresse source $a$\\
$o_3(a)$ & Troisième octet de l'adresse $a$\\
$N_Z$, $N_R$ & Nombre de zones et nombre d'enregistrements publiés\\
$B$, $Q$ & Durée du bail DHCP et durée de quarantaine d'une adresse\\
\basfilet
\end{tabular}
\end{table}
